\documentclass[12pt]{article}
\usepackage{amsmath,amssymb,amsfonts}
\usepackage[margin=1in]{geometry}

\begin{document}

\begin{center}
{\large\bfseries 3rd Irish Mathematical Olympiad\\
5 May 1990}
\end{center}

\bigskip
\noindent\textbf{Paper 1}
\bigskip

\begin{enumerate}
\item Given a natural number $n$, calculate the number of rectangles in the plane, the coordinates of whose vertices are integers in the range 0 to $n$, and whose sides are parallel to the axes.

\item A sequence of primes $a_n$ is defined as follows: $a_1 = 2$, and, for all $n \ge 2$, $a_n$ is the largest prime divisor of $a_1a_2\cdots a_{n-1} + 1$. Prove that $a_n \ne 5$ for all $n$.

\item Determine whether there exists a function $f : \mathbb{N} \to \mathbb{N}$ (where $\mathbb{N}$ is the set of natural numbers) such that \[f(n) = f(f(n-1)) + f(f(n+1)),\] for all natural numbers $n \ge 2$.

\item The real number $x$ satisfies all the inequalities \[2^k < x^k + x^{k+1} < 2^{k+1}\] for $k = 1, 2, \ldots, n$. What is the greatest possible value of $n$?

\item Let $ABC$ be a right-angled triangle with right-angle at $A$. Let $X$ be the foot of the perpendicular from $A$ to $BC$, and $Y$ the mid-point of $XC$. Let $AB$ be extended to $D$ so that $|AB| = |BD|$. Prove that $DX$ is perpendicular to $AY$.

\end{enumerate}

\newpage

\begin{center}
{\large\bfseries 3rd Irish Mathematical Olympiad\\
5 May 1990}
\end{center}

\bigskip
\noindent\textbf{Paper 2}
\bigskip

\begin{enumerate}
\setcounter{enumi}{5}
\item Let $n$ be a natural number, and suppose that the equation \[x_1x_2 + x_2x_3 + x_3x_4 + x_4x_5 + \cdots + x_{n-1}x_n + x_nx_1 = 0\] has a solution with all the $x_i$'s equal to $\pm 1$. Prove that $n$ is divisible by 4.

\item Let $n \ge 3$ be a natural number. Prove that \[\frac{1}{3^3} + \frac{1}{4^3} + \cdots + \frac{1}{n^3} < \frac{1}{12}.\]

\item Suppose that $p_1 < p_2 < \ldots < p_{15}$ are prime numbers in arithmetic progression, with common difference $d$. Prove that $d$ is divisible by 2, 3, 5, 7, 11 and 13.

\item Let $t$ be a real number, and let \[a_n = 2\cos\left(\frac{t}{2^n}\right) - 1, \quad n = 1, 2, 3, \ldots\] Let $b_n$ be the product $a_1a_2a_3\cdots a_n$. Find a formula for $b_n$ that does not involve a product of $n$ terms, and deduce that \[\lim_{n\to\infty} b_n = \frac{2\cos t + 1}{3}.\]

\item Let $n = 2k - 1$, where $k \ge 6$ is an integer. Let $T$ be the set of all $n$-tuples \[\mathbf{x} = (x_1, x_2, \ldots, x_n), \text{ where, for } i = 1, 2, \ldots, n, \; x_i \text{ is 0 or 1.}\] For $\mathbf{x} = (x_1, \ldots, x_n)$ and $\mathbf{y} = (y_1, \ldots, y_n)$ in $T$, let $d(\mathbf{x}, \mathbf{y})$ denote the number of integers $j$ with $1 \le j \le n$ such that $x_j \ne y_j$. (In particular, $d(\mathbf{x}, \mathbf{x}) = 0$). Suppose that there exists a subset $S$ of $T$ with $2^k$ elements which has the following property: given any element $\mathbf{x}$ in $T$, there is a unique $\mathbf{y}$ in $S$ with $d(\mathbf{x}, \mathbf{y}) \le 3$. Prove that $n = 23$.

\end{enumerate}

\end{document}
